\documentclass{article}
\usepackage[UTF8]{ctex}
\usepackage{amsmath,amssymb,amsfonts}
\usepackage{enumitem}
\usepackage{geometry}
\geometry{a4paper, margin=1in}
% === 错题本排版优化（不参与题目标识） ===
\setlength{\parindent}{0pt}
\setlength{\parskip}{0.35em plus 0.10em minus 0.08em}
\setlength{\emergencystretch}{3em}
\setlist{itemsep=0.25em, topsep=0.40em, parsep=0pt}
\allowdisplaybreaks[2]
% === 错题本排版优化结束 ===
\begin{document}
\section*{西安电子科技大学2025年高等代数真题}
\section*{181．西安电子科技大学2025年高等代数试卷}

一．填空题．每题5分，共25分．

\subsection{题目1}
1．实系数多项式 \(f(x)=x^{3}+p x+q\) 有一个虚根 \(3+2 \mathrm{i}\) ，另外两个根为 \(\_\_\_\_\)．

\subsection{题目2}
2．\(A=\left(\begin{array}{llll}0 & 0 & 4 & 1 \\ 0 & 0 & 1 & 4 \\ 4 & 1 & 0 & 0 \\ 1 & 4 & 0 & 0\end{array}\right), A\) 的特征值为 \(\_\_\_\_\)

\subsection{题目3}
3．\(A, B\) 为 3 阶矩阵，\(|A|=2,|B|=3\) ，则 \(\left|\begin{array}{cc}O & 2 A \\ B & O\end{array}\right|=\)

\subsection{题目4}
4．\(A\) 为4阶实对称矩阵，满足 \(A^{2}+3 A=O\) ，且 \(r(A)=2\) ，则 \(A\) 在复数域 \(\mathbb{C}\) 上的合同规范形为
\(\_\_\_\_\)，在 \(\mathbb{R}\) 上的合同规范形为 \(\_\_\_\_\)．

\subsection{题目5}
5．\(A\) 为 \(n\) 阶矩阵，\(r(A)=r\) ，且 \(A^{2}=A\) ，则 \(\operatorname{tr}(A)=\) \(\_\_\_\_\)．

\subsection{题目6}
二．（10 分）计算 \(n(n \geq 2)\) 阶行列式
\[
D_{n}=\left|\begin{array}{cccccc}
1+x^{2} & x & & & & \\
x+x^{2} & 1+x^{2} & x & & & \\
& x & 1+x^{2} & x & & \\
& & \ddots & & \ddots & \\
& & & x & 1+x^{2} & x \\
& & & & x & 1+x^{2}
\end{array}\right| .
\]

\subsection{题目7}
三．（15 分）已知 \(\alpha_{1}=(4,1,1,1,0)^{\prime}, \alpha_{2}=(8,2,4,3,0)^{\prime}, \alpha_{3}=(4,1,5,3,0)^{\prime}\) 为齐次线性方程组
\[
\left\{\begin{array}{l}
x_{1}-3 x_{2}+x_{3}-2 x_{4}-x_{5}=0 \\
-3 x_{1}+9 x_{2}-3 x_{3}+6 x_{4}+4 x_{5}=0 \\
2 x_{1}-6 x_{2}+2 x_{3}-4 x_{4}-2 x_{5}=0
\end{array}\right.
\]
的解向量．

（1） 试判断该方程组的解是否都可以由 \(\alpha_{1}, \alpha_{2}, \alpha_{3}\) 线性表示？试说明理由。
（2） 给出该方程组的一个基础解系，尽可能多地含 \(\alpha_{1}, \alpha_{2}, \alpha_{3}\) 。

\subsection{题目8}
四．（10 分）解答如下问题：

（1） \(A^{*}\) 为 \(A\) 的伴随矩阵，\(A\) 为 \(n\) 阶矩阵，证明：
\[
r\left(A^{*}\right)= \begin{cases}n, & r(A)=n \\ 1, & r(A)=n-1 \\ 0, & r(A)<n-1\end{cases}
\]
（2） \(A=\left(a_{i j}\right)_{3 \times 3}\) 是非零矩阵，\(A_{i j}\) 为 \(a_{i j}\) 的代数余子式，且满足 \(A_{i j}+a_{i j}=0, i, j=1,2,3\) ，求 \(|A|\) ．

\subsection{题目9}
五。（10 分）\(B\) 为 \(n\) 阶实正定矩阵，\(C\) 为 \(n \times k\) 实矩阵且列向量无关，记 \(A=\left(\begin{array}{cc}B & C \\ C^{\prime} & O\end{array}\right)\) ，证明：二次型\(f(X)=X^{\prime} A X\) 正负惯性指数分别为 \(n, k\) 。

\subsection{题目10}
六．（10 分）\(a_{1}, a_{2}, \cdots, a_{n}\) 为 \(n\) 个互不相同的整数，证明：\(f(x)=\prod_{i=1}^{n}\left(x-a_{i}\right)-1\) 在有理数域上不可约．

\subsection{题目11}
七．（10 分）在线性空间 \(\mathbb{R}^{2 \times 2}\) 中，记
\[
V_{1}=\left\{\left(\begin{array}{cc}
a & -a \\
b & c
\end{array}\right): a, b, c \in \mathbb{R}\right\}, V_{2}=\left\{\left(\begin{array}{cc}
a & b \\
-a & c
\end{array}\right): a, b, c \in \mathbb{R}\right\} .
\]
证明：

（1） \(V_{1}, V_{2}\) 都是 \(\mathbb{R}^{2 \times 2}\) 的子空间。
（2） 分别求 \(V_{1}+V_{2}, V_{1} \cap V_{2}\) 的一组基和维数。

\subsection{题目12}
八．（15 分）设 \(V\) 为数域 \(P\) 上的 \(n\) 维线性空间，\(\sigma, \tau\) 为 \(V\) 上的线性变换，且 \(\sigma\) 在数域 \(P\) 上有 \(n\) 个互异的特征值。证明：

（1） \(\sigma\) 的特征向量均是 \(\tau\) 的特征向量的充要条件为 \(\sigma \tau=\tau \sigma\) ．
（2） 若 \(\sigma \tau=\tau \sigma\) ，则 \(\tau\) 为 \(I, \sigma, \sigma^{2}, \cdots, \sigma^{n-1}\) 的线性组合。

\subsection{题目13}
九．（15 分）\(\sigma\) 为 \(n\) 维复线性空间 \(V\) 上的线性变换，且存在非零特征值．证明：存在 \(\sigma\) 的不变子空间\(V_{1}, V_{2}\) ，
使得 \(V=V_{1} \oplus V_{2}\) ，且 \(\sigma\) 在 \(V_{1}\) 上的限制 \(\sigma \mid V_{1}\) 是可逆的，在 \(V_{2}\) 上的限制 \(\sigma \mid V_{2}\) 是幂零的．

\subsection{题目14}
十．（15 分）已知矩阵 \(\left(\begin{array}{rrrr}1 & 0 & 0 & 0 \\ 0 & a & a & 0 \\ a-2 & 0 & 1 & 0 \\ 0 & 1 & 0 & 0\end{array}\right)\) 可对角化，求 \(a\) 的值．

\subsection{题目15}
十一．（15 分）\(A\) 为 3 阶实矩阵，\(A^{\prime} A=A A^{\prime}, A \neq A^{\prime}\) 。证明：存在正交矩阵 \(P\) ，使得
\[
P^{\prime} A P=\left(\begin{array}{ccc}
a & 0 & 0 \\
0 & b & c \\
0 & -c & b
\end{array}\right), a, b, c \text { 为实数. }
\]
\end{document}
